\documentclass[12pt]{amsart}

\usepackage{lmodern}
\usepackage{amsmath,amssymb,amsthm}
\usepackage{hyperref}
\usepackage{enumitem}

\theoremstyle{definition}
\newtheorem{definition}{Definition}[section]
\newtheorem{example}[definition]{Example}
\theoremstyle{remark}
\newtheorem{remark}[definition]{Remark}
\theoremstyle{plain}
\newtheorem{theorem}[definition]{Theorem}
\newtheorem{lemma}[definition]{Lemma}
\newtheorem{corollary}[definition]{Corollary}
\newtheorem{proposition}[definition]{Proposition}

%% Main-result environments
\newtheorem{mainthm}{Theorem}
\makeatletter
\AtBeginDocument{%
  \renewcommand{\themainthm}{\Alph{mainthm}}%
}
\makeatother

\newcommand{\re}{\operatorname{Re}}
\newcommand{\im}{\operatorname{Im}}
\newcommand{\CC}{\mathbb{C}}
\newcommand{\RR}{\mathbb{R}}
\newcommand{\NN}{\mathbb{N}}
\newcommand{\LP}{\mathrm{LP}}

\tolerance=1500
\emergencystretch=2em

\begin{document}

\title[Convergence Boundary for Real-Zero Entire Functions]{%
  Finite-Evidence Records and the Convergence Boundary\\
  for Real-Zero Entire Functions}

\author{Lin Tao}
\date{2026}

\begin{abstract}
Let $L$ be a Laguerre--P\'olya entire function of order one with a given real,
simple, square-summable zero set $\{\pm\lambda_n\}_{n\geq 1}$.
We establish two complementary theorems about sequences of entire functions
matching a finite evidence record derived from $L$.

\textbf{Theorem~\ref{thm:neg}} (per-$N$ non-identifiability): for every $N\geq 1$,
the finite evidence record $\mathcal{E}_N$ (matching the first $k_N$ zeros and
the first $J_N$ even Taylor coefficients of $L$) does not pin $L$ down.
Concretely, there exists an entire function $F$ satisfying $\mathcal{E}_N$ with
$\sup_{|z|\leq R_N}|F(z)-L(z)|\geq\varepsilon_N > 0$, where $R_N\geq 2\lambda_{k_N+1}$.
The construction is explicit: an implicit-function-theorem branch on the tail
Hadamard product, matched by a scaled Vandermonde system, separated by a
Cauchy coefficient estimate on the first unmatched log-power-sum.

\textbf{Theorem~\ref{thm:pos}} (conditional convergence): if a sequence
$(F_N)_{N\geq 1}$ satisfying $\mathcal{E}_N$ additionally has a local uniform
bound (H-bound) and a tail-reciprocal-square control (H-tail), then $F_N\to L$
locally uniformly, and the zeros of $F_N$ converge to those of $L$ by
Hurwitz's theorem.

We apply this framework to the Connes--Consani--Moscovici (CCM)
spectral-triple determinant sequence~\cite{CCM2025}, identify the three
conditions that are open in~\cite{CCM2025}, and record that specializing
$L=\Xi$ (the shifted Riemann xi function) is exactly equivalent to RH and
is never assumed. None of these results prove, disprove, or claim progress
toward the Riemann Hypothesis.
\end{abstract}

\maketitle
\tableofcontents

%---------------------------------------------------------------------
\section{Introduction}
\label{sec:intro}
%---------------------------------------------------------------------

A Laguerre--P\'olya entire function of order one,
\begin{equation}
  \label{eq:LP}
  L(z) = C\prod_{n=1}^{\infty}\!\left(1-\frac{z^2}{\lambda_n^2}\right),
  \quad C>0,\quad 0<\lambda_1<\lambda_2<\cdots,\quad
  \sum_{n=1}^{\infty}\frac{1}{\lambda_n^2}<\infty,
\end{equation}
is the canonical form of an entire function whose zeros are all real and whose
Hadamard genus is one.  The class includes many functions arising in spectral
theory and analytic number theory, most famously $\Xi(z)=\xi(1/2+iz)$
--- but only \emph{conditionally}: $L=\Xi$ requires all zeros of $\Xi$ to be
real, which is the Riemann Hypothesis (RH).  We never assume this.  RH is
$[\mathrm{OUT}]$ in this paper.

A natural strategy for studying $L$ is to approximate it by entire functions
that agree with $L$ on a growing finite record --- the first $k_N$ zeros and
the first $J_N$ even Taylor coefficients.  Such \emph{finite-evidence-record
sequences} arise naturally: the CCM determinant sequence $\det_{\mathrm{reg}}
(\mathfrak{D}_{\lambda,N}-z)$~\cite{CCM2025} satisfies conditions of this
type for $\Xi$ (conditional on RH).

\medskip
\noindent\textbf{The central question.}
Does the finite record $\mathcal{E}_N$ force a sequence to converge to $L$?

The present paper answers this in two directions:

\begin{itemize}
\item \textbf{Negative (Theorem~\ref{thm:neg}).}  No: for each fixed $N$,
  the record $\mathcal{E}_N$ alone leaves a Hadamard tail degree of freedom
  large enough to push a witness function $\varepsilon_N$ away from $L$ on a
  disk $|z|\leq R_N$ that already contains all $k_N$ matched zeros.

\item \textbf{Positive (Theorem~\ref{thm:pos}).}  Yes, provided the sequence
  additionally satisfies a local uniform bound and a tail-reciprocal-square
  control: under those two hypotheses Montel's theorem gives a normal family,
  and Hurwitz's theorem then transfers real-zero location.
\end{itemize}

\medskip
\noindent\textbf{Relation to the CCM program.}
Connes, Consani, and Moscovici~\cite{CCM2025} establish that the regularized
determinant of their spectral-triple Dirac operator satisfies
$\det_{\mathrm{reg}}(\mathfrak{D}_{\lambda,N}-z)=-i\lambda^{-iz}\hat{\xi}(z)$,
where $\hat{\xi}$ is the Fourier transform of $\xi$ (entire, all zeros real =
spectrum).  The \emph{open step} in~\cite{CCM2025} is that a suitably
normalized sequence of such determinants converges to $\Xi$.  The
$\lambda^{-iz}$ phase factor preserves zeros but not the locally uniform
limit; ``suitably normalized'' is load-bearing.  Section~\ref{sec:ccm}
identifies (H-bound), (H-uorder), and (H-div) as the three missing
conditions, and explains why (H-div) over $\Xi$ specifically would imply RH.

\medskip
\noindent\textbf{Scope and limitations.}  This paper works with an
\emph{abstract} Laguerre--P\'olya target $L$; no property of the Riemann
zeta function is assumed.  Specializing $L=\Xi$ is RH-conditional and is
always labelled as such.  The Suzuki meromorphic target $z^2\xi/\xi'$
(poles at the zeros of $\xi'$)~\cite{Suzuki2026} is out of scope: Hurwitz's
theorem does not apply to meromorphic limits, and a separate pole/residue
version of the positive theorem would be needed.

%---------------------------------------------------------------------
\section{Setting}
\label{sec:setting}
%---------------------------------------------------------------------

Throughout, $L$ denotes a fixed Laguerre--P\'olya function as in~\eqref{eq:LP}.
Write $a_n:=\lambda_n^{-2}$ for the reciprocal-square parameters; the
hypothesis $\sum a_n<\infty$ is the Hadamard summability condition.

\begin{definition}[Finite evidence record]
\label{def:record}
Fix integers $k_N, J_N\geq 1$.  An entire function $F$ \emph{satisfies the
finite evidence record} $\mathcal{E}_N$ (relative to $L$) if:
\begin{enumerate}[label=(\arabic*)]
\item $F$ is entire of order one.
\item $F$ is even: $F(-z)=F(z)$.
\item $F$ is real on the real axis: $F(\bar{z})=\overline{F(z)}$.
\item All zeros of $F$ are real.
\item The first $k_N$ positive zeros of $F$ (ordered by size) equal
  $\lambda_1\leq\lambda_2\leq\cdots\leq\lambda_{k_N}$.
\item Base-point normalization: $F(0)=L(0)=C$.
\item Even Taylor coefficients match:
  $F^{(2j)}(0)=L^{(2j)}(0)$ for $j=0,1,\ldots,J_N$.
\end{enumerate}
\end{definition}

\begin{remark}
\label{rem:log-power-sums}
By the Hadamard factorization, conditions~(1)--(4) force $F$ to have the form
$F(z)=C_F\prod_{n}(1-z^2/r_n^2)$ for some real sequence $(r_n)$.  The
Taylor conditions~(6)--(7) are equivalent to matching the \emph{log-power-sums}
$P_r(F):=\sum_n r_n^{-2r}=\sum_n a_n^r$ for $r=1,\ldots,J_N$ (together with
$F(0)=C$).  This is the key reformulation used in the proof.
\end{remark}

%---------------------------------------------------------------------
\section{Theorem A: per-\texorpdfstring{$N$}{N} non-identifiability}
\label{sec:neg}
%---------------------------------------------------------------------

\begin{mainthm}[Finite record does not identify $L$]
\label{thm:neg}
For every $N\geq 1$ there exist $\varepsilon_N>0$, $R_N\geq 2\lambda_{k_N+1}$,
and an entire function $F$ satisfying $\mathcal{E}_N$ with
\[
  \sup_{|z|\leq R_N}\lvert F(z)-L(z)\rvert \;\geq\; \varepsilon_N.
\]
Equivalently, the fiber $\{F : F \text{ satisfies } \mathcal{E}_N\}$ is not
contained in any $\varepsilon$-neighbourhood of $L$ on $\{|z|\leq R_N\}$.
\end{mainthm}

\begin{remark}[Per-$N$, not sequence divergence]
\label{rem:per-N}
Theorem~\ref{thm:neg} is a \emph{per-$N$} statement: for each fixed $N$ there
is a single witness $F$ (depending on $N$) separated from $L$.  It does not
assert that any particular \emph{sequence} $(F_N)$ fails to converge locally
uniformly.  Under the extra assumption $k_N\to\infty$, the radii $R_N$ escape
to infinity, but the witnesses are not required to form a sequence satisfying
all $\mathcal{E}_N$ simultaneously.  The content is precisely that
$\mathcal{E}_N$ \emph{alone} leaves a tail degree of freedom.
\end{remark}

\begin{proof}
Fix $N$; write $k:=k_N$, $J:=J_N$, $a_m:=\lambda_{k+m}^{-2}$ for $m\geq 1$.
We construct the witness in three steps.

\medskip
\textbf{Step~1 (one-parameter tail deformation).}
For $m>J$ define
\[
  b_m(c) \;:=\; a_m\bigl(1+c/m\bigr)^{-2},\qquad c\in(-1,1).
\]
At $c=0$, $b_m(0)=a_m$.  Note that $\sum_m b_m(c)<\infty$ for all $c$ in a
neighbourhood of $0$ (by dominated convergence, since $\sum a_m<\infty$).

\medskip
\textbf{Step~2 (IFT matching).}
Let $u=(u_1,\ldots,u_J)$ be the reciprocal squares of the first $J$ free
tail zeros, and define the log-power-sum matching system
\[
  \Phi_r(u,c) \;:=\; \sum_{\ell=1}^{J} u_\ell^r + \sum_{m>J} b_m(c)^r
                \;-\; \sum_{m\geq 1} a_m^r \;=\; 0,\qquad r=1,\ldots,J.
\]
At $(u^{(0)},c)=((a_1,\ldots,a_J),0)$ we have $\Phi=0$.  The Jacobian
$(\partial\Phi_r/\partial u_\ell)|_{u^{(0)},0} = r\,a_\ell^{r-1}$ is an
exact scaled Vandermonde matrix with nodes $a_1>\cdots>a_J>0$ (all distinct
since $\lambda_n$ are simple), so its determinant is nonzero.  By the implicit
function theorem there exist $\delta>0$ and a $C^1$ branch $u(c)$ with
$\Phi(u(c),c)=0$ for $0<c<\delta$.

Define
\[
  F_c(z) \;:=\; C\prod_{n=1}^{k}\!\left(1-\frac{z^2}{\lambda_n^2}\right)
               \cdot\prod_{\ell=1}^{J}\!\bigl(1-u_\ell(c)\,z^2\bigr)
               \cdot\prod_{m>J}\!\bigl(1-b_m(c)\,z^2\bigr).
\]
By construction $F_c$ is entire of order one, even, real on $\RR$, has all
real zeros, first $k$ positive zeros equal to $\lambda_1,\ldots,\lambda_k$,
and satisfies $F_c(0)=C$ and $P_r(F_c)=P_r(L)$ for $r=1,\ldots,J$.  Thus
$F_c$ satisfies $\mathcal{E}_N$.

\medskip
\textbf{Step~3 (Cauchy separation).}
The first unmatched log-power-sum is
$\Delta_{J+1}(c):=P_{J+1}(F_c)-P_{J+1}(L)$.  Differentiating the IFT
solution at $c=0$ gives
\[
  \Delta_{J+1}'(0) \;=\;
  -(J+1)\sum_{m>J} \frac{2a_m}{m}\,\prod_{\ell=1}^{J}(a_m-a_\ell),
\]
where every factor $\prod_{\ell}(a_m-a_\ell)$ has sign $(-1)^J$ (since
$a_m<a_J<\cdots<a_1$), so the sum is nonzero (all terms have the same sign).
Shrinking $\delta$ we have $\Delta_{J+1}(c)\neq 0$ for $0<c<\delta$.

Since $F_c^{(2J+2)}(0)\neq L^{(2J+2)}(0)$, the difference $F_c-L$ is an
entire function with a zero of order exactly $2J+2$ at the origin.  Writing
$F_c(z)-L(z)=A_c z^{2J+2}+O(z^{2J+4})$ with $A_c\neq 0$, Cauchy's
inequality on the disk $|z|\leq R$ gives
\[
  \sup_{|z|\leq R}\lvert F_c(z)-L(z)\rvert \;\geq\; |A_c|\,R^{2J+2}.
\]
Setting $R_N:=\max\{2\lambda_{k+1},\,(\varepsilon/|A_c|)^{1/(2J+2)}\}$ for any
fixed $\varepsilon>0$ completes the proof.
\end{proof}

%---------------------------------------------------------------------
\section{Theorem B: conditional convergence}
\label{sec:pos}
%---------------------------------------------------------------------

\begin{definition}[Hypotheses for convergence]
\label{def:hypotheses}
Let $(F_N)_{N\geq 1}$ be a sequence of entire functions satisfying
$\mathcal{E}_N$.  We say the sequence satisfies:
\begin{itemize}
\item[(H-norm)] \emph{Base-point convergence}: $F_N(z_0)\to L(z_0)\neq 0$
  for some $z_0\in\CC$.
\item[(H-bound)] \emph{Local uniform boundedness}: for every $R>0$ there
  exists $M_R>0$ such that $\sup_N\sup_{|z|\leq R}|F_N(z)|\leq M_R$.
\item[(H-tail)] \emph{Tail reciprocal-square control}: writing
  $(r_{n,N})_{n\geq 1}$ for the positive zeros of $F_N$ and setting
  $\tilde{a}_{n,N}:=r_{n,N}^{-2}$, we have
  $\sum_{n=1}^{\infty}|\tilde{a}_{n,N}-a_n|<C$ uniformly in $N$ for some
  constant $C<\infty$.
\end{itemize}
\end{definition}

\begin{mainthm}[Conditional convergence to $L$]
\label{thm:pos}
Let $(F_N)_{N\geq 1}$ satisfy $\mathcal{E}_N$ and hypotheses
\textup{(H-norm)}, \textup{(H-bound)}, and \textup{(H-tail)}.
Then $F_N\to L$ locally uniformly on $\CC$, and the zeros of $F_N$
converge to the zeros of $L$ \textup{(}with multiplicity\textup{)} by
Hurwitz's theorem.
\end{mainthm}

\begin{proof}
\textbf{Normal family.}  By~(H-bound), the family $\{F_N\}$ is locally
uniformly bounded; Montel's theorem gives a locally uniformly convergent
subsequence $F_{N_j}\to G$ for some entire function $G$.

\textbf{Identification $G=L$.}  The function $G$ is entire of order at most
one (by~(H-tail) and the Hadamard factorization: $G=C_G\prod_n(1-z^2/r_n^2)$
where $r_n$ are real, and $\sum r_n^{-2}=\lim \sum\tilde{a}_{n,N}<\infty$
by~(H-tail)).  Moreover $G$ is even and real on $\RR$ (by uniform convergence
from~(2)--(3) of $\mathcal{E}_N$).  By~(H-norm), $G(z_0)=L(z_0)\neq 0$,
so $G\not\equiv 0$.  The matching conditions~(5)--(7) of $\mathcal{E}_N$
(finite zero agreement + Taylor coefficients) pass to the limit: $G$ agrees
with $L$ on $k_N\to\infty$ zeros and on all even Taylor coefficients, so
by the Hadamard factorization uniqueness theorem (an even entire function of
order one is determined by its zero set and one nonzero value) we have $G=L$.

Since every subsequence of $(F_N)$ has a further subsequence converging to
$L$, the full sequence $F_N\to L$ locally uniformly.

\textbf{Zeros.}  $L$ has only simple zeros (the $\lambda_n$ are distinct
by assumption), so Hurwitz's theorem applies: for each $\lambda_n$ and
$\varepsilon>0$, all zeros of $F_N$ in $|\lambda_n-z|<\varepsilon$ converge
to $\lambda_n$ for large $N$.
\end{proof}

%---------------------------------------------------------------------
\section{Application to the CCM determinant sequence}
\label{sec:ccm}
%---------------------------------------------------------------------

Connes, Consani, and Moscovici~\cite{CCM2025} prove the determinant identity
\begin{equation}
  \label{eq:ccm}
  \det_{\mathrm{reg}}(\mathfrak{D}_{\lambda,N}-z)
  \;=\; -i\,\lambda^{-iz}\,\hat{\xi}(z),
\end{equation}
where $\hat{\xi}$ is the Fourier transform of $\xi$ and is an entire function
with all zeros real (= spectrum of $\mathfrak{D}_{\lambda,N}$).
The target for convergence is $\Xi(z)=\xi(1/2+iz)$, not $\hat{\xi}$.
The $\lambda^{-iz}$ phase factor in~\eqref{eq:ccm} preserves the zero set
but alters the locally uniform limit: ``suitably normalized'' in~\cite{CCM2025}
is the load-bearing step that Theorems~\ref{thm:neg}--\ref{thm:pos} analyze.

\begin{table}[h]
\centering
\caption{Status of convergence conditions for the CCM sequence
  (assuming RH for~(5) and~(H-div); all else stated in~\cite{CCM2025}).}
\label{tab:ccm}
\begin{tabular}{lll}
\hline
Condition & Status & Reference \\
\hline
(1) Entire & \textsc{proved} & CCM det.\ identity \\
(2)--(3) Even, real on $\RR$ & \textsc{proved} & Symmetry of $\mathfrak{D}$ \\
(4) Real zeros & \textsc{proved} & CCM Thm., real spectrum \\
(5) First $k_N$ zeros agree & \textsc{open} & Numerical; no analytic proof \\
(6) Base-point normalization & \textsc{partial} & $\lambda^{-iz}$ phase obstacle \\
(H-norm) $F_N(z_0)\to\Xi(z_0)$ & \textsc{open} & The ``suitably normalized'' step \\
(H-bound) Local uniform bound & \textsc{open} & No tail envelope in~\cite{CCM2025} \\
(H-tail) Tail reciprocal-square & \textsc{open} & No spectral tail bound \\
\hline
\end{tabular}
\end{table}

The three open conditions (H-norm), (H-bound), (H-tail) are precisely
what Theorem~\ref{thm:pos} needs to conclude $F_N\to\Xi$ locally uniformly.

\begin{remark}[Why (H-div) over $\Xi$ implies RH]
\label{rem:hdiv-rh}
In the original formulation~\cite{CCM2025}, a condition (H-div) required
the zero divisors of $F_N$ to converge to the zero divisor of $\Xi$.
This forces every zero of $\Xi$ to be a real limit of real zeros of $F_N$.
If the zero set of $\Xi$ is \emph{complete} (i.e., $\Xi$ has no zeros off
the real line), then this is exactly the statement that all zeros of $\Xi$
are real --- which is RH.  Thus (H-div) over $\Xi$ is an RH-strong hypothesis.
We do not impose it.  RH is $[\mathrm{OUT}]$.
\end{remark}

\begin{remark}[Suzuki meromorphic target out of scope]
\label{rem:suzuki}
Suzuki~\cite{Suzuki2026} proves the implication: if $W(a,\theta;z)\to
z^2\xi(1/2-iz)/\xi'(1/2-iz)$ locally uniformly on compact subsets, then RH
follows (Corollary~6).  The target $z^2\xi/\xi'$ is \emph{meromorphic} (poles
at the zeros of $\xi'$).  Theorem~\ref{thm:pos} applies only to entire targets;
the Hurwitz step fails for meromorphic limits because zeros of approximants
can accumulate at poles rather than at zeros of the target.  A separate
pole/residue version of the positive theorem is needed for the Suzuki target.
\end{remark}

%---------------------------------------------------------------------
\section{Limitations}
\label{sec:limitations}
%---------------------------------------------------------------------

\begin{enumerate}[label=(\alph*)]
\item \textbf{$\Xi$-specific version requires RH.}  Specializing $L=\Xi$
  in Theorem~\ref{thm:neg} requires $\Xi$ to be an LP function of the
  form~\eqref{eq:LP} --- i.e., all zeros of $\Xi$ are real --- which is RH.
  The abstract-$L$ statement is RH-free; the $\Xi$ instance is not claimed.

\item \textbf{Simplicity of zeros.}  The IFT step (Step~2 of the proof of
  Theorem~\ref{thm:neg}) uses the scaled Vandermonde nonsingularity, which
  requires the $\lambda_n$ to be simple (distinct).  For $\Xi$, zero
  simplicity is an open problem; it is an explicit hypothesis on $L$.

\item \textbf{Suzuki target excluded.}  See Remark~\ref{rem:suzuki}.

\item \textbf{Sequence convergence not ruled out.}  Theorem~\ref{thm:neg}
  is a per-$N$ statement (Remark~\ref{rem:per-N}).  It does not rule out
  locally uniform convergence $F_N\to L$ for a particular sequence satisfying
  all $\mathcal{E}_N$ simultaneously; the CCM sequence is believed to converge
  after suitable normalization.
\end{enumerate}

%---------------------------------------------------------------------
\section{Open problems}
\label{sec:open}
%---------------------------------------------------------------------

\begin{enumerate}
\item \textbf{LP functions with dependent zero sets.}
  The proof of Theorem~\ref{thm:neg} uses the simplicity and summability
  of $\{\lambda_n\}$.  Can the result be extended to LP functions with
  double zeros, or to functions of order greater than one?

\item \textbf{Effective $R_N$.}  The proof gives $R_N\geq 2\lambda_{k_N+1}$;
  can one give an explicit $R_N$ as a function of $N$, $\lambda_1,\ldots,
  \lambda_{k_N+1}$, and $J_N$, with a certified $\varepsilon_N>0$?

\item \textbf{Uniform tail bound from spectral theory.}
  For the CCM sequence, can (H-tail) be derived from spectral estimates
  on $\mathfrak{D}_{\lambda,N}$ (e.g., via Suzuki's $\lambda(a)$ bounds
  from~\cite{Suzuki2026})?  Closing this gap would unconditionally supply
  the tail-reciprocal-square control, reducing Theorem~\ref{thm:pos} to
  requiring only (H-norm) and (H-bound) for the CCM case.
\end{enumerate}

{\sloppy
\begin{thebibliography}{99}

\bibitem{CCM2025}
A.~Connes, C.~Consani, and H.~Moscovici,
\emph{Zeta spectral triples},
arXiv:2511.22755, 2025.

\bibitem{CvS2025}
A.~Connes and W.~van Suijlekom,
\emph{Spectral truncations in noncommutative geometry and operator systems},
Commun.\ Math.\ Phys.~\textbf{406} (2025), article 15.

\bibitem{Suzuki2026}
M.~Suzuki,
\emph{Weil's quadratic form via the screw function},
arXiv:2606.09096, 2026.

\bibitem{Levin1996}
B.~Ya.~Levin,
\emph{Lectures on Entire Functions},
Translations of Mathematical Monographs, vol.~150,
Amer.\ Math.\ Soc., Providence, RI, 1996.

\bibitem{Conway1978}
J.~B.~Conway,
\emph{Functions of One Complex Variable},
2nd ed., Springer, New York, 1978.

\end{thebibliography}
}

\end{document}
